\section{Local coding obstructions}\label{sec:local-coding}

The construction programme of~\cite{FirstPaper2026} needs coding stages that
transmit fewer bits than the uncoded data.  To test a proposed stage, one must
specify what each operation reads, what each receiver knows, and what
input-dependent information selects the operations.  This section gives three
obstructions: a separator bound for adaptive queries of original rows; a
distance bound for fixed coordinate-XOR broadcasts after bijective row
recoding; and a helper-or-deficit bound for two ordered nonlinear stages.

Each result concerns its stated finite coding model.  None asserts that every
fast tape transposer has that form.  In particular, the first model reads
original rows, not stored rows of coded workspace.

Throughout this section, $d,b\in\mathbb N$, $[d]=\{0,\ldots,d-1\}$,
\[
  \mathcal X=\{0,1\}^{[d]\times[d]\times[b]},\qquad P=d^2b.
\]
A row belongs to $Q=\{0,1\}^{[d]\times[b]}$, and
$(x^{\mathsf T})_{ijt}=x_{jit}$.  The finite-coordinate adapter in
\texttt{MatrixView} uses the transpose and encoding of the first paper; it
does not replace their specification.  Except where positivity is stated,
zero dimensions and zero entry widths are permitted.

For a probability law $w$ on a finite correct-input family $F\subseteq\mathcal X$,
write $X\sim w$ and
\[
  H(X)=-\sum_{x\in F}w(x)\log_2 w(x),\qquad D_w=P-H(X),
\]
with the usual zero-mass convention.  The weights may vanish and need not be
independent across rows, entries, or bits.  Under the uniform law on nonempty
$F$, write $D=P-\log_2|F|$.  All entropies in this section are in bits; the Lean
entropy functional uses natural logarithms and is divided by $\log 2$ at the
endpoints.

\subsection{Adaptive queries of original rows}\label{subsec:local-adaptive}

\begin{definition}[Public original-row strategy]\label{def:local-public-strategy}
A local query consists of a support $A\subseteq[d]$ and a function
$q:\mathcal X\to\{0,1\}$ such that
\[
  x|_A=y|_A\quad\Longrightarrow\quad q(x)=q(y).
\]
The support may have any finite size.  A public strategy assigns, to every
finite binary history $h$, either a local query or a halt instruction.  When
it selects $q$, the next transmitted bit is $q(x)$ and the history becomes
$h\,q(x)$.  An execution is complete only if the strategy then halts; the
halt instruction adds no bit.

A helper $z:\mathcal X\to\mathcal Z$, with $\mathcal Z$ finite, selects a
strategy $s_{z(x)}$.  Receiver $j$ has a decoding function of precisely
$z(x)$, its own original row $x_j$, and the completed transcript $M(x)$.
For a declared family $F$, the correctness requirement is that, for every
$x\in F$, the selected strategy terminates with transcript $M(x)$ and every
receiver returns $(x^{\mathsf T})_j$.  There is no termination or correctness
requirement outside $F$.
\end{definition}
\formal{ProtocolTruncation}{TransposeFollowup.PublicStrategy,TransposeFollowup.PublicStrategy.Executes}
\formal{LocalProtocol}{TransposeFollowup.LocalQuery}

Thus the helper and public history determine the next support and the
stopping decision.  Any further input-dependent control must be conveyed by
query answers.  Query evaluation and decoding are unrestricted; the cost is
the number of transmitted bits, including those used for control.

For a completed transcript $m$ under helper value $z$, let $G(z,m)$ be the
simple graph on $[d]$ obtained by completing each queried support to a clique.
Only the queries actually encountered along $m$ contribute edges.  In Lean,
this graph is obtained from the strategy truncated at $|m|$; the execution
hypothesis ensures that no query of the completed run is lost.  A separator
selection is any function $S(z,m)\subseteq[d]$.  No balance condition is
assumed.  If $\mathcal C(G-S)$ denotes the components of the induced graph
outside $S$, define
\begin{equation}\label{eq:local-separator-loss}
  \mathcal L(G,S)=db|S|+b\sum_{C\in\mathcal C(G-S)}|C|^2.
\end{equation}
For random input $X$, abbreviate $Z=z(X)$, $G_X=G(Z,M(X))$, and
$S_X=S(Z,M(X))$.

\begin{theorem}[Adaptive separator bound]\label{thm:local-adaptive-separator}
For the strategy of Definition~\ref{def:local-public-strategy}, correct on
$F$, and every probability law $w$ on $F$,
\begin{equation}\label{eq:local-adaptive-separator}
  \mathbb E_w|M(X)|
  \ \ge\ P-\mathbb E_w\mathcal L(G_X,S_X)-2D_w-2H(Z).
\end{equation}
The separator may depend arbitrarily on the helper and completed transcript.
Correctness is required on all of the declared family $F$, including its
zero-probability inputs.  For uniform $X$ on nonempty $F$, the same statement
holds with $D_w=D$.
\end{theorem}
\formal{GraphSeparatorBound}{TransposeFollowup.graph_expected_separator_bound_of_distribution,TransposeFollowup.graph_expected_separator_bound}

The bound excludes a fixed fractional saving whenever expected separator loss, input
deficit, and helper entropy are all $o(P)$.  It therefore tests a proposed
coding stage without imposing linearity on its queries or decoders.

\begin{proofoutline}
First replace the possibly infinite strategies by finite protocols, truncating
all of them at the maximum transcript length attained on $F$.  Finiteness of
$F$ makes this maximum finite, and the truncation preserves every specified
execution, its cost, and its queried supports.  This step imposes no assumption
on the original strategies outside $F$.

Fix a helper value $z$, transcript $m$, and separator assignment $u$ attained
on $F$.  Let $R$ be the \emph{virtual full fiber}: all
inputs producing $m$ in the fixed finite protocol for $z$ and having separator
assignment $u$.  Membership in $R$ does not require the input's actual helper
to be $z$, nor does it require correctness.  Every query along $m$ meets at
most one component of $G(z,m)-S(z,m)$ outside the separator.  Consequently,
component assignments taken from different elements of $R$ can be combined
freely while keeping $u$ fixed.  The locality equation for each query shows,
one step at a time, that the combined input follows the same transcript.
Thus $R$ is exactly the product of its component projections, with the
separator coordinates fixed.

Let $A$ be a nonempty correct subfamily with this helper, transcript, and
separator assignment.  No product structure is asserted for $A$.  On $A$,
each component's original rows determine the columns requested by its
receivers.  Fixed separator rows likewise determine the separator receivers'
columns, so those rows and columns have zero entropy.

For the remaining coordinates, each component's owned-and-demanded set
covers an intercomponent bit twice and an internal bit once.  Adding one
copy of each internal bit gives a twofold entropy cover.  Decoding reduces
each owned-and-demanded entropy to the owned-row entropy.  The product
structure of $R$ bounds the sum of these entropies by $\log_2|R|$; the
internal bits contribute at most $b\sum_C|C|^2$.  Under the uniform law on
$A$, this gives
\begin{equation}\label{eq:local-leaf-count}
  2\log_2|A|\le \log_2|R|+
        b\sum_{C\in\mathcal C(G-S)}|C|^2.
\end{equation}
The product property is used only for $R$, not for the correlated family $A$.

For the aggregation, append the separator rows, in a fixed order, to each
transcript.  A leaf label $a$ consists of the helper and this augmented
transcript.  Write $f_a$ for its correct-fiber size, $V_a$ for its virtual
full-fiber size, $q_a=|m|+db|S|$ for the augmented length, and
$r_a=b\sum_C|C|^2$ for the internal loss.  For each fixed helper $z$, the
virtual fibers are disjoint, and the augmented transcripts form a prefix-free
code.  The latter follows because completed public transcripts are
prefix-free and, once $m$ is known, the appended row-string length is fixed.
Hence
\[
  \sum_{a:\,z(a)=z} V_a\le 2^P,
  \qquad
  \sum_{a:\,z(a)=z}2^{-q_a}\le 1.
\]
Let $A_*$ be the random leaf label.  Conditional Gibbs inequalities for these
two budgets and the finite-fiber entropy inequality give
\begin{align*}
 H(X)&\le H(A_*)+\mathbb E\log_2 f_{A_*},\\
 H(A_*)&\le H(Z)+P-\mathbb E\log_2 V_{A_*},\\
 H(A_*)&\le H(Z)+\mathbb E q_{A_*}.
\end{align*}
Adding the last two inequalities and using twice the first together with
\eqref{eq:local-leaf-count} yields
$2H(X)\le P+2H(Z)+\mathbb E(q_{A_*}+r_{A_*})$.
Since $q_{A_*}+r_{A_*}=|M(X)|+\mathcal L(G_X,S_X)$, rearrangement proves
\eqref{eq:local-adaptive-separator}.  All averaging uses the supplied law;
zero-mass fibers contribute zero.
\end{proofoutline}

\begin{corollary}[Worst-case cost and attained helper range]\label{cor:local-separator-worst}
Under the same strategy assumptions, let $F\ne\varnothing$ and let
$B,L,g\in\mathbb N$.  If every $x\in F$ satisfies
$|M(x)|\le B$ and $\mathcal L(G_x,S_x)\le L$, and
$|z(F)|\le 2^g$, then
\[
  B\ge P-L-2(D+g).
\]
Only helper values attained on $F$ are charged; the ambient helper alphabet
is still required to be finite.
\end{corollary}
\formal{GraphSeparatorBound}{TransposeFollowup.graph_worst_case_separator_bound}
This form gives a direct rejection test from worst-case bounds on a candidate's
broadcast, separator loss, and attained helper range.
\begin{proofoutline}
Use the uniform law in Theorem~\ref{thm:local-adaptive-separator}, bound both
expectations by their pointwise bounds, and use $H(Z)\le\log_2|z(F)|\le g$.
\end{proofoutline}

\begin{theorem}[Separate sharpness of the two information charges]\label{thm:local-separator-sharpness}
In the stated original-row model, the coefficient of $D$ and the coefficient
of $H(Z)$ in \eqref{eq:local-adaptive-separator} cannot separately be replaced
by a constant smaller than $2$, with the other terms unchanged.  Each assertion
has a zero-transcript witness with $d=2$, empty separator, and every $b>0$.
\end{theorem}
\formal{SeparatorSharpness}{TransposeFollowup.SeparatorSharpness.deficit_coefficient_at_least_two,TransposeFollowup.SeparatorSharpness.helper_coefficient_at_least_two}
Thus a stronger obstruction must use further structure or change other terms
of the bound; reducing either information coefficient alone is impossible.
\begin{proofoutline}
The empty interaction graph has two singleton components, so its loss is
$2b$, while $P=4b$.  For the deficit coefficient, take the family
$x_{01}=x_{10}$, of cardinality $2^{3b}$, and a constant helper.  Each receiver
already has its column in its own row.  The transcript has length zero,
$D=b$, and a coefficient $c$ would require $0\ge 2b-cb$.
For the helper coefficient, take the full input cube and the helper
$Z=x_{01}\oplus x_{10}$.  Each receiver recovers the other off-diagonal
entry from its own row and $Z$, again with no transcript.  Here $D=0$ and,
under the uniform input law, $H(Z)=b$.  The same inequality forces $c\ge2$.
\end{proofoutline}

\subsubsection{A treewidth consequence}

Use ordinary finite tree decompositions, with the natural-valued convention
$\operatorname{tw}(\varnothing)=0$: $\operatorname{tw}(G)$ is the least
$t\in\mathbb N$ for which a tree decomposition has bags of size at most
$t+1$.  The decomposition tree has arbitrary finite shape.

\begin{lemma}[Integral fragmentation]\label{lem:local-treewidth-fragmentation}
For every finite simple graph $G$ on $d$ vertices and every $L\in\mathbb N$,
there is $S\subseteq V(G)$ such that
\[
  (L+1)|S|\le (\operatorname{tw}(G)+1)d,
  \qquad |C|\le L\quad(C\in\mathcal C(G-S)).
\]
For a graph on $[d]$, a separator can consequently be chosen so that
\[
  (L+1)\mathcal L(G,S)
       \le bd\bigl((\operatorname{tw}(G)+1)d+L(L+1)\bigr).
\]
\end{lemma}
\formal{Treewidth}{TransposeFollowup.treewidth_fragmentation,TransposeFollowup.treewidth_separator_loss}
\begin{proofoutline}
Assign each graph vertex to one bag that contains it, giving integer weights
of total $d$ to the decomposition tree.  The weighted-tree fragmentation
lemma removes a set of tree vertices whose number, multiplied by $L+1$, is
at most $d$, leaving weight at most $L$ in each remaining tree component.
Remove their bags from $G$.  The running-intersection and edge-cover
properties transfer the component bound to $G$, while each removed bag costs
at most $\operatorname{tw}(G)+1$ vertices.  Finally,
$\sum_C|C|^2\le L\sum_C|C|\le Ld$ gives the loss estimate.  This is a
standard tree-decomposition argument; its role here is to provide the exact
integral budget used in the next optimization.
\end{proofoutline}

\begin{corollary}[Savings force large average treewidth]\label{cor:local-treewidth}
Assume $d,b>0$ and the strategy hypotheses of
Theorem~\ref{thm:local-adaptive-separator}.  For any probability law $w$ on
$F$, suppose
$\mathbb E_w|M(X)|\le(1-\varepsilon)P$, where $\varepsilon\in\mathbb R$.
Put
\[
  \eta=\frac{D_w+H(Z)}{P},\qquad
  a=\max\{\varepsilon-2\eta,0\}.
\]
Then
\begin{equation}\label{eq:local-treewidth}
  \mathbb E_w\operatorname{tw}(G_X)
       \ge \frac d4 a^2+\frac a2-1.
\end{equation}
\end{corollary}
\formal{AdaptiveTreewidth}{TransposeFollowup.average_treewidth_lower_bound_of_distribution,TransposeFollowup.average_treewidth_lower_bound}
\begin{proofoutline}
For each integer $L\ge0$, choose the separator supplied by
Lemma~\ref{lem:local-treewidth-fragmentation}.  It is a function of the public
leaf graph, so it is an allowed separator.  Average its loss bound and apply
the separator theorem.  With $E=\mathbb E\operatorname{tw}(G_X)$ and
$a_0=\varepsilon-2\eta$, the resulting family of tests is
\[
  d(L+1)a_0\le d(E+1)+L(L+1)\qquad(L\in\mathbb N).
\]
If $a_0\le0$, nonnegativity of treewidth suffices.  Otherwise set
$r=da_0/2$ and $L=\lfloor r\rfloor$.  The elementary identity
\[
  (L+1)(2r-L)-(r^2+r)=(r-L)(L+1-r)\ge0
\]
gives $d(E+1)\ge r^2+r$ from the test at this $L$, which is
\eqref{eq:local-treewidth}.
\end{proofoutline}

As $d\to\infty$, with $b>0$ and $D_w+H(Z)=o(P)$, a fixed fractional saving
therefore requires average treewidth $\Omega(d)$.  Alternatively, separators
of size $o(d)$ whose remaining components all have size $o(d)$ have loss
$o(P)$, provided these estimates hold uniformly over the completed runs.
Such original-row strategies cannot yield a fixed fractional saving.
Graphs without such separators remain candidates for this model.

The first paper's construction discussion asks for locally encoded index
coding and suggests, heuristically, that hyperfinite interaction graphs offer
only negligible savings~\cite[section ``The construction side'']{FirstPaper2026}.
The separator and treewidth bounds make this restriction quantitative for
nonlinear queries, public adaptivity, helpers, and correlated correct sets.
They leave the coding problem on bounded-degree expanders open and do not
strengthen Theorem~P or prove the Tape Coding Conjecture.

\subsection{Bijective row recoding and fixed XOR packets}\label{subsec:local-xor}

The preceding separator argument leaves highly connected row graphs open.
For a more restrictive broadcast format, an independent obstruction applies
even when the row graph has bounded degree and large treewidth.

\begin{definition}[Recoded fixed-XOR code]\label{def:local-recoded-xor}
For each $i\in[d]$, choose an arbitrary bijection $A_i:Q\to Q$ and write
$y_i=A_i(x_i)$.  Fix $B\in\mathbb N$ packets, each either one coordinate
$y_{ijt}$ or the XOR of two coordinates of $y$.  Packet choices are fixed
before the input is supplied.  For a simple graph $G$ on $[d]$, locality means
that the row indices of a pair packet are equal or adjacent in $G$; singleton
packets impose no additional edge condition.  Receiver $j$ receives the
complete $B$-bit broadcast and its own \emph{original} row $x_j$, and must
recover $(x^{\mathsf T})_j$ for every $x\in\mathcal X$.
\end{definition}
\formal{RecodedXor}{TransposeFollowup.RecodedXorCode,TransposeFollowup.RecodedXorCode.Correct,TransposeFollowup.RecodedXorCode.LocalTo}
\formal{XorPackets}{TransposeFollowup.XorPacket,TransposeFollowup.XorSystem}

The bijections and the decoders need not be linear.  Repeated packets,
same-row pairs, and self-XORs are allowed.  A self-XOR is a permitted constant
zero packet.  There is no helper in this model, and correctness is on the
entire input cube, not merely a high-entropy subfamily.

For $\ell\in\mathbb N$, call $(i,j)$ \emph{$\ell$-near} if there exists a
walk in $G$ from $i$ to $j$ of length strictly below $\ell$, and let
\[
  Q_\ell=b\,\bigl|\{(i,j)\in[d]^2:(i,j)\text{ is $\ell$-near}\}\bigr|.
\]
This walk-based definition excludes disconnected pairs; it does not assign
them distance zero.  In particular $Q_0=0$.

\begin{theorem}[Distance bound after arbitrary row recoding]\label{thm:local-xor-distance}
Every correct code of Definition~\ref{def:local-recoded-xor} that is local to
$G$ satisfies, for all $\ell\in\mathbb N$,
\begin{equation}\label{eq:local-xor-distance}
  (\ell+1)P\le(\ell+1)B+P+\ell Q_\ell,
  \qquad
  B\ge P-\frac{P+\ell Q_\ell}{\ell+1}.
\end{equation}
No positivity or connectedness hypothesis is needed for the integral form.
\end{theorem}
\formal{XorComponentCount}{TransposeFollowup.RecodedXorCode.broadcast_lower_bound_scaled,TransposeFollowup.RecodedXorCode.broadcast_lower_bound}

Small graph neighborhoods limit the savings even after nonlinear row
recoding.  The next two corollaries turn this bound into explicit restrictions
on bounded-degree graphs and on paths and cycles.

\begin{proofoutline}
Form the variable graph on the $P$ transformed bit coordinates: distinct
endpoints of a pair packet are adjacent.  A coordinate appearing as a
singleton packet is an anchor.  Let $u$ be the number of components containing
no anchor.  The broadcast together with one transformed bit from each
unanchored component determines every transformed bit: equal broadcasts make
the difference of two assignments constant on each component, and anchors or
chosen representative bits determine that constant.  Injective finite counting
therefore gives $P\le B+u$.

An unanchored component can be flipped simultaneously without changing the
broadcast.  Pull this flip back through the row bijections.  If its represented
row set is $R$, original rows outside $R$ do not change.  Their receivers have
the same side information and broadcast before and after the flip, so all
columns indexed outside $R$ also stay fixed.  The raw change is therefore
confined to $R\times R\times[b]$.  This holds at every basepoint, despite
nonlinearity of the row bijections.

If the variable component contains at most $\ell$ coordinates, any two of
its coordinates are joined by a simple path of length below $\ell$.
Projection to row indices collapses same-row edges and maps other edges to
edges of $G$.  Thus its raw changes are confined to the $Q_\ell$ near bits.
Let $u_{\mathrm{small}}$ count such unanchored components.  Their disjoint
nonempty flips generate $2^{u_{\mathrm{small}}}$ distinct transformed
assignments and hence the same number of raw assignments.  Because each flip
fixes every far bit at every basepoint, their compositions do so as well.
The orbit fits into $Q_\ell$ binary coordinates, giving
$u_{\mathrm{small}}\le Q_\ell$.

Disjointness and nonemptiness of all unanchored components imply
\[
  (\ell+1)u\le P+\ell u_{\mathrm{small}}.
\]
Indeed, a large component pays at least $\ell+1$ coordinates, while a small
one pays at least one and receives the additional allowance $\ell$.
Combining these inequalities proves \eqref{eq:local-xor-distance}.
\end{proofoutline}

\begin{corollary}[Bounded degree]\label{cor:local-xor-degree}
For a correct code local to $G$, if $d\ge2$ and $G$ has maximum degree at
most $\Delta\in\mathbb N$, then
\[
 B\ge P-\frac{6\log(\max\{2,\Delta\})}{\log d}\,P.
\]
The logarithm base is immaterial to this ratio.  For fixed $\Delta$ and
$b>0$, no code in this model achieves a fixed positive fractional saving as
$d$ grows.
\end{corollary}
\formal{XorDegreeBound}{TransposeFollowup.RecodedXorCode.fixed_degree_broadcast_lower_bound}
\begin{proofoutline}
Put $\delta=\max\{2,\Delta\}$ and
$k=\lfloor\lfloor\log_\delta d\rfloor/2\rfloor$.  Counting walks bounds each
radius-$k$ ball by $2\delta^k$, so $Q_{k+1}\le2bd\delta^k$.
The integer estimate $(k+1)\delta^k\le d$ yields
$(k+2)P\le(k+2)B+3P$ from the distance theorem.
Finally $\log d\le2(k+2)\log\delta$ converts this to the displayed bound.
The argument is uniform in $b$ and does not assume regularity or
connectedness of $G$.
\end{proofoutline}

\begin{corollary}[Paths, cycles, and their embedded subgraphs]\label{cor:local-xor-cycle}
For a correct code local to $G$, assume $d>0$.  If there is an injective graph
homomorphism from $G$ into a cycle on $m\ge1$ vertices, then
\begin{equation}\label{eq:local-xor-cycle}
  B\ge P-bd\bigl(\sqrt{8(d+3)}-4\bigr).
\end{equation}
In particular this holds for paths, cycles, and their subgraphs, with arbitrary
vertex labels.  The same bound holds whenever every $\ell$-near ball has at
most $2\ell-1$ vertices for every integer $\ell\ge1$.
\end{corollary}
\formal{XorPathCycleBound}{TransposeFollowup.RecodedXorCode.linear_ball_broadcast_lower_bound,TransposeFollowup.RecodedXorCode.cycle_embedding_broadcast_lower_bound,TransposeFollowup.RecodedXorCode.subcycle_broadcast_lower_bound,TransposeFollowup.RecodedXorCode.cycle_broadcast_lower_bound,TransposeFollowup.RecodedXorCode.path_broadcast_lower_bound}
For these graphs, the maximum possible saving is $O(P/\sqrt d)$, improving
the $O(P/\log d)$ error supplied by the bounded-degree estimate.
\begin{proofoutline}
An injective graph homomorphism sends a near ball into the corresponding
cycle ball.  Walking fewer than $\ell$ steps in a cycle gives at most
$2\ell-1$ possible integer offsets; collisions can only reduce this number.
Hence $Q_k\le bd(2k-1)$ for every $k\ge1$.  The distance theorem gives
\[
  P-B\le bd\,\frac{d+k(2k-1)}{k+1}.
\]
Choose $k=\lfloor\sqrt{\lfloor(d+3)/2\rfloor}\rfloor$, put $j=k+1$, and
let $t=\sqrt{(d+3)/2}$, where the quotient defining $t$ is real.
Then $k\le t<j$, $j\ge2$, and
\[
 \frac{d+k(2k-1)}{k+1}
    =4t-5+\frac{2(j-t)^2}{j}\le4t-4
    =\sqrt{8(d+3)}-4.
\]
The estimate includes $d=1$.
\end{proofoutline}

The bounded-degree conclusion applies even to expanders, but only for the
specified packet format.  It does not answer the first paper's question about
\emph{arbitrary} linear functionals supported on two rows: one such functional
may involve many coordinates, and a single common row bijection need not
convert a whole family of them into pairwise coordinate XORs.
The first paper's XOR-frame model concerns changes of address frame in a
gadget, a different restriction.  The bounded-degree corollary rules out a
fixed fractional saving through independent bijective row preprocessing
followed by a fixed coordinate-XOR broadcast local to a bounded-degree graph.

\subsection{Two nonlinear ordered stages and a global helper}\label{subsec:local-two-stage}

A different construction strategy lets later stages act on coded rows.
The following model allows two coupled stages on one bank of $d$ rows.
Its obstruction is independent of runtime: all allowed maps are granted as
primitives.

\begin{definition}[Triangular and causal row maps]\label{def:local-triangular-causal}
For a finite row alphabet $Q$ and a permutation $\sigma$ of $[d]$, a triangular
row map has the form
\[
  T(u)_i=\phi_{i,u_{<i}}(u_i),\qquad
  u_{<i}=(u_j)_{\sigma(j)<\sigma(i)},
\]
where $\phi_{i,v}:Q\to Q$ is a bijection for every fixed earlier-row tuple
$v$.  Dependence on $v$ is arbitrary.  A causal row map in an order $\tau$ has
the form
\[
  C(p)_j=\gamma_j\bigl((p_l)_{\tau(l)\le\tau(j)}\bigr),
\]
with arbitrary functions $\gamma_j$; no injectivity is required.
These formulas refer to the incoming row tuple of each stage.
\end{definition}
\formal{TriangularRows}{TransposeFollowup.TriangularRows,TransposeFollowup.TriangularRows.apply,TransposeFollowup.CausalRows,TransposeFollowup.CausalRows.apply}

\begin{definition}[Two-stage branch and helper selection]\label{def:local-two-stage}
A two-stage branch consists of three row permutations
$\pi_{\mathrm{in}},\pi_{\mathrm{mid}},\pi_{\mathrm{out}}$, arbitrary local
maps $a_i,c_i:Q\to Q$, a triangular map $T$, local bijections
$\lambda_i:Q\to Q$, and a causal map $C$.  It acts by
\begin{align*}
 u_i&=a_i(x_{\pi_{\mathrm{in}}(i)}),&
 v_i&=T(u)_i,& w_i&=\lambda_i(v_i),\\
 p_j&=w_{\pi_{\mathrm{mid}}(j)},&
 q_j&=C(p)_j,& \Phi(x)_i&=c_i(q_{\pi_{\mathrm{out}}(i)}).
\end{align*}
A helper-selected architecture is a function $z:\mathcal X\to\mathcal Z$
and a branch $\Phi_z$ for each $z\in\mathcal Z$.  Here $\mathcal Z$ may be
infinite.  For $Q=\{0,1\}^{[d]\times[b]}$, correctness on $F$ means
$\Phi_{z(x)}(x)=x^{\mathsf T}$ for every $x\in F$.
\end{definition}
\formal{TwoStageRows}{TransposeFollowup.TwoStageRows,TransposeFollowup.TwoStageRows.apply}
\formal{TwoStageHelper}{TransposeFollowup.TwoStageHelper,TransposeFollowup.TwoStageHelper.CorrectOn}

All maps and both stage orders may depend on the single helper value.
Once that value is fixed, they are fixed functions of the displayed inputs.
The first stage is required to be bijective in its own row argument, and the
middle local maps are bijections.  The input maps, output maps, and second
stage may discard information.  An extra bank, a global operation between the
stages, or another input-dependent controller is not part of this definition.

\begin{lemma}[Block injectivity]\label{lem:local-block-injective}
A triangular map $T$ is block-injective: for every $I\subseteq[d]$,
\[
  u|_{I^c}=u'|_{I^c},\qquad T(u)|_I=T(u')|_I
       \quad\Longrightarrow\quad u=u'.
\]
It is a bijection on the finite row space.  Postcomposition by rowwise
bijections preserves block injectivity, and the inverse of a block-injective
bijection is block-injective.  In particular, the inverse of
$W(u)_i=\lambda_i(T(u)_i)$ has this property.
\end{lemma}
\formal{TriangularRows}{TransposeFollowup.TriangularRows.block_injective,TransposeFollowup.TriangularRows.equiv,TransposeFollowup.blockInjective_rowwise,TransposeFollowup.blockInjective_inverse}
\formal{TwoStageRows}{TransposeFollowup.TwoStageRows.inverse_block_injective}
This supplies the recovery property used in the next counting lemma directly
from the allowed row maps, without imposing it as an extra hypothesis.
\begin{proofoutline}
Induct on the triangular order.  Outside $I$ the two incoming rows already
agree.  Inside $I$, earlier rows agree by induction, so the same bijection
$\phi_{i,u_{<i}}$ is applied to both own-row inputs; equality of the outputs
forces equality of those inputs.  Taking $I=[d]$ gives global injectivity,
and finiteness gives bijectivity.  Rowwise postcomposition is cancelled one
row at a time.  For the inverse statement, apply block injectivity of the
forward map with $I^c$ in place of $I$.
\end{proofoutline}

\begin{lemma}[A record for one branch]\label{lem:local-branch-count}
If one fixed two-stage branch is correct on an arbitrary family
$F\subseteq\mathcal X$, then for every integer $0\le r\le d$,
\begin{equation}\label{eq:local-branch-count}
  |F|\le 2^{(d-r)db+r^2b}=2^{P-br(d-r)}.
\end{equation}
The family may be empty or correlated.
\end{lemma}
\formal{TwoStageRows}{TransposeFollowup.TwoStageRows.correct_family_card}
This bounds how many inputs one branch can handle; partitioning by the helper
will then bound both its required range and its entropy.
\begin{proofoutline}
Let $\tau$ be the second-stage order.  Its last $r$ incoming rows determine
sets
\begin{align*}
 I&=\{i:\tau(\pi_{\mathrm{mid}}^{-1}(i))\ge d-r\},\\
 R&=\pi_{\mathrm{in}}(I),\qquad
 J=\{i:\tau(\pi_{\mathrm{out}}(i))\ge d-r\}.
\end{align*}
All three sets have size $r$.  The input rows $u_I$ depend only on the original
rows $x_R$.  By causality, the final output rows outside $J$ depend only on
$w_{I^c}$.

Record $w_{I^c}$ and the raw rectangle $x_{R\times J\times[b]}$.  There are
$2^{(d-r)db+r^2b}$ possible records.  To prove injectivity on $F$, suppose two
correct inputs have the same record.  Equality of $w_{I^c}$ gives equality of
their final output rows outside $J$, hence equality of their raw columns
outside $J$.  The recorded rectangle supplies the remaining entries of the
original rows in $R$, so those complete rows agree.  Therefore $u_I$ agrees.
Block injectivity of $W^{-1}$, together with equality of $w_{I^c}$, now gives
equality of all of $w$.  The two complete outputs coincide; correctness and
injectivity of transpose imply that the original inputs coincide.

The support sets come from the specified permutations and causal stage, as
implemented in \texttt{FixedSupportSandwich}.  The exponent identity uses
$r\le d$.
\end{proofoutline}

\begin{theorem}[Exact helper-or-deficit toll]\label{thm:local-helper-count}
For a helper-selected two-stage architecture correct on $F\subseteq\mathcal X$
and every integer $0\le r\le d$,
\begin{equation}\label{eq:local-helper-count}
  |F|\,2^{br(d-r)}\le |z(F)|\,2^P.
\end{equation}
In particular,
\[
  |F|\,2^{b\lfloor d^2/4\rfloor}\le |z(F)|\,2^P.
\]
If $D\in\mathbb R$, $g\in\mathbb N$, $|F|\ge2^{P-D}$, and
$|z(F)|\le2^g$, then
\begin{equation}\label{eq:local-helper-deficit}
  D+g\ge b\lfloor d^2/4\rfloor.
\end{equation}
The counting statements include $F=\varnothing$; the density premise in the
last assertion implies $F\ne\varnothing$.  Only the finite attained helper
image is counted.
\end{theorem}
\formal{TwoStageHelper}{TransposeFollowup.TwoStageHelper.correct_family_count_toll,TransposeFollowup.TwoStageHelper.quarter_count_toll,TransposeFollowup.TwoStageHelper.quarter_helper_bound_of_density}
A small helper range therefore forces a large correct-family deficit.
For example, a branch architecture correct on the full cube needs at least
$2^{b\lfloor d^2/4\rfloor}$ attained helper values.
\begin{proofoutline}
Partition $F$ by its helper value.  On each part the architecture is one fixed
branch, so Lemma~\ref{lem:local-branch-count} bounds its cardinality by
$2^{P-br(d-r)}$.  Sum over the attained helper values and multiply by
$2^{br(d-r)}$.  Choosing $r=\lfloor d/2\rfloor$ uses the integer identity
$r(d-r)=\lfloor d^2/4\rfloor$.  The density conclusion follows by taking
logarithms of positive quantities.  No rounding of the real deficit $D$ is
needed, and no factorization of $F$ is used.
\end{proofoutline}

\begin{theorem}[Actual helper entropy under arbitrary input laws]\label{thm:local-helper-entropy}
Under the same architecture and correctness hypotheses, let $X$ have any
probability law on $F$ and put $Z=z(X)$ and $D_w=P-H(X)$.  Then, for every
integer $0\le r\le d$,
\begin{equation}\label{eq:local-helper-entropy}
  D_w+H(Z)\ge br(d-r),
  \qquad D_w+H(Z)\ge b\lfloor d^2/4\rfloor.
\end{equation}
The ambient helper alphabet may be infinite.  Its entropy is computed on the
finite attained subtype.  Uniform input gives the corresponding inequalities
with $D_w=P-\log_2|F|$.
\end{theorem}
\formal{TwoStageHelperEntropy}{TransposeFollowup.TwoStageHelper.helper_entropy_deficit_bound_of_distribution,TransposeFollowup.TwoStageHelper.quarter_helper_entropy_deficit_bound_of_distribution,TransposeFollowup.TwoStageHelper.helper_entropy_deficit_bound,TransposeFollowup.TwoStageHelper.quarter_helper_entropy_deficit_bound}
\begin{proofoutline}
Every fiber of the attained-helper map on $F$ has size at most
$C=2^{P-br(d-r)}$ by the branch-count lemma.  The shared finite-fiber entropy
inequality gives $H(X)\le H(Z)+\log_2 C$.  In Lean this inequality follows
from conditional Gibbs with comparison weight $C^{-1}$ on each input of a
helper fiber; each such fiber has comparison mass at most one.  It is valid
for arbitrary nonnegative weights, so specializing to a normalized law does
not require positive mass at every input.  Rearrangement gives the first
inequality, and $r=\lfloor d/2\rfloor$ gives the second.
\end{proofoutline}

For $d\to\infty$, $b>0$, and $D_w=o(P)$, the entropy bound requires
$H(Z)\ge(1/4-o(1))P$.  A construction in this model must therefore supply
that much helper information even if its helper values are very unevenly
distributed.  The theorem does not cover more stages, extra banks, or an
unrestricted causal network.

The first paper identifies a global oracle hint in the plain-versus-robust
gap, and its proposed recursion needs savings at both encoder and
decoder~\cite[sections ``Plain versus robust'' and ``The construction
side'']{FirstPaper2026}.  The two-stage results quantify a helper-or-deficit
toll for one explicit architecture; they neither remove the general hint nor
strengthen the native pass lower bound.  The three coding models in this
section impose different restrictions.  A construction escaping one
obstruction must still be checked against the hypotheses of the others.
